\section{A coupled three-level counterexample to uniform sharpness}
The saturation in~\eqref{eq:sharp} is special. A metric can become badly
conditioned near a defective eigenvalue that is \emph{nonzero}; in this
situation, inversion need not approach singularity. The next lower bound
relates metric conditioning to an ordinary eigenvector overlap in any
dimension.

\begin{lemma}[Overlap obstruction for every admissible metric]
\label{lem:overlapn}
Let $H$ be diagonalizable with real spectrum and let $v,w$ be
Euclidean-unit right eigenvectors for distinct eigenvalues. Write
$c=|v^\dagger w|<1$. For every $\eta>0$ satisfying
$H^\dagger\eta=\eta H$,
\begin{equation}\label{eq:overlapn}
 \kappa_2(\eta)\geq\frac{1+c}{1-c}.
\end{equation}
\end{lemma}
\begin{proof}
Distinct real eigenvalues give $v^\dagger\eta w=0$ by the
intertwining equation. Restrict the positive quadratic form of $\eta$
to $S=\mathrm{span}\{v,w\}$. In an orthonormal basis of $S$, its
matrix $\eta_S$ obeys the two-vector inequality
$\kappa_2(\eta_S)\geq(1+c)/(1-c)$ proved in
Proposition~\ref{prop:overlap}. The Rayleigh quotient bound
$\lambda_{\min}(\eta)\leq\lambda_{\min}(\eta_S)
\leq\lambda_{\max}(\eta_S)\leq\lambda_{\max}(\eta)$
then gives $\kappa_2(\eta)\geq\kappa_2(\eta_S)$.
\end{proof}

\begin{theorem}[Collapse of the optimized certificate without loss of invertibility]
\label{thm:counterexample}
For $0<\varepsilon\leq1/2$, consider the coupled three-mode chain
\begin{equation}\label{eq:threelevel}
 A_\varepsilon=
 \begin{pmatrix}
 1&1&0\\
 0&1+\varepsilon&1\\
 0&0&2
 \end{pmatrix}.
\end{equation}
It is diagonalizable with distinct, positive eigenvalues and has a
positive admissible metric. Its minimal spectral modulus is $m=1$, but
\begin{align}
 \kappa_*(A_\varepsilon)&\geq
 \frac{\sqrt{1+\varepsilon^2}+1}
 {\sqrt{1+\varepsilon^2}-1},\label{eq:threecond}\\
 0<\sup_{\eta>0:\,A_\varepsilon^\dagger\eta
 =\eta A_\varepsilon}
 \frac{m}{\sqrt{\kappa_2(\eta)}}
 &\leq\frac{\varepsilon}
 {\sqrt{1+\varepsilon^2}+1}
 \longrightarrow0,\label{eq:threecert}\\
 r_{\mathrm{sing}}(A_\varepsilon)
 &\geq\frac{2}{\sqrt{15}}>0.
 \label{eq:threeradius}
\end{align}
Consequently the ratio of the best metric-based certificate to the
exact Euclidean singularity radius tends to zero. The endpoint
$A_0$ is defective at eigenvalue $1$, but is itself invertible.
\end{theorem}
\begin{proof}
The eigenvalues are $1$, $1+\varepsilon$, and $2$; they are
positive and distinct on the stated interval. Eigenvectors for the
first two eigenvalues may be chosen as
\begin{equation}
 v=(1,0,0)^\mathsf{T},\qquad
 w=\frac{(1,\varepsilon,0)^\mathsf{T}}
 {\sqrt{1+\varepsilon^2}},\qquad
 c=|v^\dagger w|=(1+\varepsilon^2)^{-1/2}.
\end{equation}
Lemma~\ref{lem:overlapn} yields~\eqref{eq:threecond}; since
$m=1$, the monotonicity of the inverse square root gives
\eqref{eq:threecert}. Direct inversion yields
\begin{equation}\label{eq:threeinverse}
 A_\varepsilon^{-1}=
 \begin{pmatrix}
 1&-u&u/2\\
 0&u&-u/2\\
 0&0&1/2
 \end{pmatrix},\qquad u=(1+\varepsilon)^{-1}\leq1.
\end{equation}
Hence
\begin{equation}
 \norm{A_\varepsilon^{-1}}_2^2
 \leq\norm{A_\varepsilon^{-1}}_F^2
 =\frac54+\frac52u^2\leq\frac{15}{4},
\end{equation}
which proves~\eqref{eq:threeradius} using~\eqref{eq:radius}.
Finally $A_0$ has determinant $2$ and a nontrivial Jordan block
at eigenvalue one.
\end{proof}

\begin{remark}[The chain is not an orthogonal direct sum]
\label{rem:coupled}
For $0<\varepsilon\leq1/2$, $A_\varepsilon$ has no common
eigenvector with $A_\varepsilon^\dagger$. Indeed, its three
right-eigendirections are spanned by
$(1,0,0)^\mathsf{T}$, $(1,\varepsilon,0)^\mathsf{T}$,
and $(1,1,1-\varepsilon)^\mathsf{T}$
(up to nonzero scaling for the last vector).
The first is not an eigenvector of $A_\varepsilon^\dagger$,
the second has a nonzero third component after application of
$A_\varepsilon^\dagger$, and the third cannot be a common
eigenvector since its nonzero first component would force the
adjoint eigenvalue to be one rather than two.
Thus no one-dimensional reducing subspace exists; in dimension three,
there is no nontrivial orthogonal block decomposition.
\end{remark}

The counterexample separates \emph{eigenvector coalescence} from
\emph{loss of invertibility}. It does not assert that all exceptional
points occur at nonzero eigenvalues or that every physical notion of
stability remains uniform. In particular, the metric itself has no
positive-definite continuation through the defective endpoint.

\subsection{Sharp asymptotic metric conditioning}
The two-eigenvector overlap bound detects quadratic divergence, but it does not
capture the full three-state geometry.  The next theorem strengthens the lower
bound and determines the exact leading constant of the optimized metric
condition number.
